\subsection{Reusing factors through a graph without four-cycles}
\label{sec:alias-compression}

The private factor for each term occurrence gives a simple sufficient
condition for coefficient separation. It is stronger than necessary.
We can reuse a private-side alias among occurrences that denote the same
logical first factor, provided a second graph satisfies an elementary
restriction. This \emph{factor graph} is different from the graph whose
divergences define the polynomial.

Put the private-side aliases in one part $A$ and the shared representatives
in another part $S$, including the common unit representative in $S$.
A nonsquare occurrence $m_i=a_is_i$ is an edge between its two factors.
Require that this bipartite graph is simple and has no cycle of length
four. In the version allowing $m_i=a_i^2$, each such square occurrence
keeps an alias that appears in no other source monomial.

\paragraph{Uniqueness of products.}
Suppose that two products of two nonsquare occurrence monomials are equal
as formal monomials. Equality determines the multisets of private-side
and shared-side factors. If either multiset is repeated, simplicity
forces the pair of edges. Otherwise the only possible ambiguity is
\[
 (as)(bt)=(at)(bs), \qquad a\ne b,\quad s\ne t.
\]
The four required edges would form a four-cycle, which is forbidden.
An isolated square occurrence is recognized by its private coordinate,
which occurs in no other source monomial. Products involving such an
occurrence are therefore uniquely recognized as well.

It follows that source-edge squares and products of distinct source
edges never collide after substitution. A source-edge/flow-edge product
is identified by its flow coordinate and its distinct source monomial.
Thus the half-incidence energy still has coefficients in
$\{-1,0,1\}$. The alias equality cycles occupy the separate quadratic
alias namespace, while the flow quadratics use flow coordinates.
The unit path supplies the only linear and constant terms. These
observations establish coefficient separation for the entire polynomial.

\paragraph{Preservation of the zero set.}
An alias is reused only among copies of the same logical normalized
variable. Those copies already agree at every zero of the uncompressed
construction. Conversely, every source zero assigns that logical value
to the reused alias. Rebuild the ordinary equality cycle on the distinct
surviving aliases, adding one dummy exactly when the class has size two.
Both maps on zero sets remain mutually inverse. All internal graph flows
are still uniquely determined by their spanning-tree equations.

\paragraph{A deterministic finite construction.}
The supplied compressor processes occurrences in their recorded order.
It keeps square occurrences separate. An occurrence with shared factor
$s$ can join a private group for the same logical variable if $s$ is new
to the group and no pair $\{s,t\}$ with an existing neighbor $t$ has
already been assigned to another private group. Among eligible groups
it chooses a largest one, with creation order breaking ties. Otherwise
it creates a new group. The neighbor-pair ledger proves inductively that
the factor graph stays simple and has no four-cycle. This is a
deterministic feasible packing, with no assertion of optimality.

On the 442-variable source with 128 hubs and square aliases, this merges
1,969 private occurrences into 1,752 groups. Rebuilding the equality
cycles removes 61 further dummy variables. The resulting complete
polynomial has 2,215 variables and 30,468 monomials, compared with
2,493 variables and 31,186 monomials before this compression. Its total
degree is four, its coefficient height is one, and its individual
variable degrees are at most four. All these counts refer to the fully
collected polynomial, independently reconstructed from the sum of
squares for twice that polynomial.

\subsection{A small example of sharing the flow witnesses}
\label{sec:two-row-example}

Let two abstract residuals be
\[
 R_1=A+B-C,\qquad R_2=D+E-F.
\]
Join their row vertices to three common hub vertices. Only two free
hub-path flows, $p$ and $q$, are needed. The half-incidence energy $\mathcal E$
is specified by
\begin{align*}
 2\mathcal E={}&R_1^2+R_2^2+(p-A-D)^2\\
              &+(q-p-B-E)^2+(C+F-q)^2.
\end{align*}
Every abstract edge symbol occurs in exactly two squared divergences.
Two different edge symbols meet in at most one divergence. Consequently
$\mathcal E$ has integer coefficients of absolute value at most one.
It is nonnegative and vanishes precisely when $R_1=R_2=0$, with the
uniquely determined flows
\[
 p=A+D,\qquad q=A+B+D+E=C+F.
\]
The last equality uses the two residual equations. Away from their
common zero set, choose the polynomial formulas
$p=A+D$ and $q=A+B+D+E$. Then the exact identity is
\[
 \mathcal E=\frac12\bigl(R_1^2+R_2^2+(R_1+R_2)^2\bigr).
\]
This illustrates why sharing hubs preserves the zero set while changing
the off-zero value identity. The general construction substitutes the
protected quadratic occurrence monomials for the six abstract symbols;
the factor-separation argument above prevents coefficient collisions
under that substitution.
